\small
\setlength{\extrarowheight}{3pt}
\renewcommand{\arraystretch}{1.3}

\begin{longtable}{
    >{\raggedright\arraybackslash}p{2.6cm}
    >{\raggedright\arraybackslash}p{5.8cm}
    >{\raggedright\arraybackslash}p{4.8cm}
}

\caption{Artigos Selecionados para a Revisão Rápida}
\label{tab:artigos_selecionados}\\

\toprule
\textbf{Citação} & \textbf{Título e Foco Principal} & \textbf{Relevância para o Projeto} \\
\midrule
\endfirsthead

\multicolumn{3}{c}{\textit{Continuação da Tabela \ref{tab:artigos_selecionados}}} \\
\toprule
\textbf{Citação} & \textbf{Título e Foco Principal} & \textbf{Relevância para o Projeto} \\
\midrule
\endhead

\midrule
\multicolumn{3}{r}{\textit{Continua na próxima página}} \\
\endfoot

\bottomrule
\endlastfoot

\rowcolor[gray]{0.96}
\cite[]{boulahia2024} &
\textit{The multi-criteria evaluation of research efforts based on ETL software} \newline
(Avaliação de múltiplos critérios (AHP) para escolha de ETLs em BI, Big Data e contexto semântico.) &
Critérios de avaliação de ETLs Open Source (M1). \\

\cite[]{nwokeji2021} &
\textit{A Systematic Literature Review on Big Data Extraction, Transformation and Loading (ETL)} \newline
(Revisão sistemática sobre ETL para Big Data, abordando abordagens e atributos de qualidade.) &
Estrutura e atributos de qualidade para ETL em Big Data (M1). \\

\rowcolor[gray]{0.96}
\cite[]{murarka2024} &
\textit{Advanced Techniques in Data Ingestion and Pipelining for Scalable Big Data Platforms} \newline
(Técnicas avançadas de ingestão e processamento para Big Data escalável, com uso de Kafka, Spark e ML/AI.) &
Demonstra a escalabilidade de ferramentas como Kafka e Spark no ecossistema Open Source (M2/M5). \\

\cite[]{septian2019} &
\textit{A Study Review of Common Big Data Architecture for Small-medium Enterprise} \newline
(Arquitetura de Big Data e pipelines para Pequenas e Médias Empresas, com foco em Open Source.) &
Justifica o foco em Open Source para PMEs (M1) e serve de base para a arquitetura (M2). \\

\rowcolor[gray]{0.96}
\cite[]{reverseetl2022} &
\textit{Reverse ETL for Improved Scalability, Observability, and Performance of Modern Operational Analytics – A Comparative Review} \newline
(Análise e importância do Reverse ETL para o fluxo inverso do Data Warehouse para sistemas operacionais.) &
Identifica a tendência emergente do Reverse ETL, essencial para a análise crítica (M5). \\


\rowcolor[gray]{0.96}
\cite[]{mbata2024survey} &
\textit{A Survey of Pipeline Tools for Data Engineering} \newline
(Levantamento abrangente das categorias ETL/ELT, ingest\~{a}o, orquestra\c{c}\~{a}o e aprendizado de m\'aquina, classificando e comparando dezenas de ferramentas de c\'odigo aberto e propriet\'arias.) &
Cat\'alogo de ferramentas por categoria (M1/M2), base para a tabela~\ref{tab:ferramentas_open_source} e para a escolha tecnol\'ogica do projeto. \\

\end{longtable}